% Section 2: METIS Seed Variance (~1,800 words)

\subsection{METIS as a Randomized Algorithm}

METIS recursive bisection \citep{karypis1998} is a heuristic graph
partitioning algorithm. Like most combinatorial optimization heuristics,
it uses randomization to escape local optima: the random seed
$s \in \{0, 1, \ldots, 2^{32}-1\}$ initializes the coarsening,
initial partitioning, and refinement phases. Different seeds can produce
partitions with different edge cuts and hence different district shapes,
even for the same input graph.

The DIA \citep{b02paper} specifies a single deterministic seed derived
from the census release identifier via SHA-256 \citep{sha256nist}:
\begin{equation}
  s_0 = \text{SHA-256}(\text{GEOID\_release\_id}) \bmod 2^{32}.
\end{equation}
This ensures reproducibility. But it raises the question: how close is
the DIA seed's output to the minimum-edge-cut partition, and how much
do output metrics vary if a different seed had been specified?

Paper B.7 \citep{b7paper} conducted the definitive empirical
characterization: 10,000 seed realizations per state, across all 50
states, for the 2020 Census congressional maps (500,000 total METIS
calls). We use the B.7 dataset directly to derive uncertainty intervals
for metric values under seed variation.

\subsection{Seed Variance in Edge Cut}

The coefficient of variation (CV) of the normalized edge cut across
10,000 seeds is the primary measure of METIS seed sensitivity. From
B.7 (Table~1):
\begin{itemize}
  \item CV $= 0.0\%$ for single-district states ($k = 1$; 7 states).
  \item Mean CV $= 0.6\%$, max $= 1.2\%$ for small states ($k \leq 5$;
    12 states).
  \item Mean CV $= 1.4\%$, max $= 4.3\%$ for medium states
    ($6 \leq k \leq 15$; 20 states).
  \item Mean CV $= 2.1\%$, max $= 3.8\%$ for large states ($k > 15$;
    11 states).
  \item CV $< 2\%$ for 48 of 50 states; Georgia ($4.3\%$) and
    North Carolina ($3.8\%$) are outliers.
\end{itemize}

Low CV in edge cut does not automatically imply low CV in downstream
metrics (PP, partisan outcomes). Edge cut is a global measure of
partition quality; PP depends on the specific geometric shape of each
district. A partition with the same edge cut as another may arrange
that cut differently, producing different district shapes.

\subsection{Bootstrap CI for Polsby-Popper}

We treat the 10,000 seed realizations as a bootstrap sample of the
seed distribution, enabling direct estimation of the sampling
distribution of any downstream metric without parametric assumptions.

For each state $s$ and seed $i$, let $PP_{s,i}$ denote the mean
Polsby-Popper score across all districts in the algorithmic plan
generated by seed $i$. The bootstrap 95\% CI for the expected mean
PP under the DIA seed is:
\begin{equation}
  \left[\hat{PP}_s - z_{0.975} \cdot \widehat{se}(PP_s),\;
        \hat{PP}_s + z_{0.975} \cdot \widehat{se}(PP_s)\right]
\end{equation}
where $\hat{PP}_s = \frac{1}{10000}\sum_{i=1}^{10000} PP_{s,i}$
and $\widehat{se}(PP_s) = \text{SD}(PP_{s,1}, \ldots, PP_{s,10000}) / \sqrt{10000}$.

\begin{proposition}
Under the assumption that METIS seed realizations are exchangeable
draws from the seed distribution, the bootstrap CI above is a valid
95\% CI for the expected PP of a uniformly random seed.
\end{proposition}

\paragraph{Bootstrap validity assumption.}
We treat the 50-state sample as exchangeable units for the purpose of
constructing a CI for the mean PP. This is appropriate if states are
drawn from a common distribution of geographic complexity --- an
assumption supported by the finding that PP variance is primarily
geographic (C.0, C.2). Three interpretations of the CI are possible:
(a)~a CI for the \emph{DIA seed's specific output} (not bootstrap ---
the DIA seed is fixed and its output is a point, not a distribution);
(b)~a CI for the \emph{expected output of a uniformly random seed},
which is what the bootstrap estimates; and (c)~a CI for a future DIA
seed in a different application. This paper reports interpretation
(b). It is the relevant quantity for legal proceedings because it
characterizes what any seed --- including the DIA seed --- is likely
to produce, and demonstrates that the DIA seed is not a cherry-picked
outlier from the seed distribution.

Table~\ref{tab:pp-seed-ci} reports 95\% bootstrap CIs for the national
mean PP and for selected states representing diverse CV regimes.

\begin{table}[htb]
\centering
\caption{95\% bootstrap confidence intervals for mean Polsby-Popper
  under seed variation. National mean is the population-weighted
  average across all 435 districts.}
\label{tab:pp-seed-ci}
\small
\begin{tabular}{lcccc}
\toprule
State / Region & Point Est.\ & 95\% CI Lower & 95\% CI Upper & CI Width \\
\midrule
\textbf{National mean} & \textbf{0.361} & \textbf{0.351} & \textbf{0.370} & 0.019 \\
\midrule
Vermont ($k=1$)      & 0.521 & 0.521 & 0.521 & 0.000 \\
Delaware ($k=1$)     & 0.498 & 0.498 & 0.498 & 0.000 \\
\midrule
Idaho ($k=2$)        & 0.486 & 0.484 & 0.488 & 0.004 \\
New Hampshire ($k=2$)& 0.463 & 0.461 & 0.466 & 0.005 \\
\midrule
Colorado ($k=8$)     & 0.432 & 0.428 & 0.436 & 0.008 \\
Maryland ($k=8$)     & 0.398 & 0.392 & 0.404 & 0.012 \\
\midrule
Florida ($k=28$)     & 0.473 & 0.468 & 0.478 & 0.010 \\
Texas ($k=38$)       & 0.459 & 0.453 & 0.465 & 0.012 \\
California ($k=52$)  & 0.421 & 0.415 & 0.427 & 0.012 \\
\midrule
Georgia ($k=14$, high CV)     & 0.419 & 0.407 & 0.431 & 0.024 \\
North Carolina ($k=14$, high CV) & 0.428 & 0.417 & 0.439 & 0.022 \\
\bottomrule
\end{tabular}
\end{table}

The national mean PP CI ($[0.351, 0.370]$) is tight: width $0.019$
or $\pm 2.6\%$ relative to the point estimate. For most states, seed
variance contributes CI widths of $0.004$--$0.012$, small relative
to the $\sim 0.10$ gap between algorithmic and enacted plans.

Georgia and North Carolina exhibit notably wider CIs ($0.022$--$0.024$)
reflecting their high edge-cut CV ($4.3\%$ and $3.8\%$). Even so,
the CI widths are small relative to the observed compactness differences
between algorithmic and enacted plans in those states.

\subsection{Bootstrap CI for Partisan Seat Share}

From B.7's partisan analysis, for the two highest-variance states:
\begin{itemize}
  \item Wisconsin ($k=8$): Democratic seat share ranges from 3/8 to 5/8
    across 10,000 seeds, with mean $3.7/8$ and SD $0.4$ seats.
    95\% CI for mean Democratic seats: $[3.0, 4.4]$ out of 8.

  \item North Carolina ($k=14$): Democratic seat share ranges from 5/14
    to 7/14 with mean $5.6/14$ and SD $0.5$ seats.
    95\% CI for mean Democratic seats: $[4.6, 6.6]$ out of 14.
\end{itemize}

The 2-seat range for these highest-variance states is an upper bound
on seed-driven partisan uncertainty across all 50 states. For 47/50
states, partisan seat-share SD is below $0.3$ at 50 seeds.

\subsection{The DIA Seed Is Not Systematically Biased}

A natural concern is whether the SHA-256-derived DIA seed $s_0$
produces systematically worse outcomes than the minimum-edge-cut seed.
From B.7, the approximation gap of the DIA seed is $2.9\% \pm 1.7\%$,
statistically indistinguishable from the median seed gap
($3.1\% \pm 1.8\%$; paired $t$-test $p = 0.41$). The DIA seed is
neither better nor worse than a random seed, as expected: SHA-256
produces outputs computationally indistinguishable from uniform random.

This means the 95\% bootstrap CI around the DIA seed's output is
centered at the same location as the CI for any other seed, and the
DIA seed is not systematically biasing any metric up or down.

\subsection{Implications for Report Formatting}

Based on the bootstrap CI analysis, we recommend the following
reporting conventions for expert witness reports and legislative
testimony:

\begin{itemize}
  \item Report Polsby-Popper as: $PP = 0.361\ (95\%\text{ CI: }[0.351, 0.370])$
    at national level, or state-level CIs from Table~\ref{tab:pp-seed-ci}.

  \item Report partisan seat share for high-CV states (GA, NC, WI) with
    explicit seed-variance CI; for all other states a single decimal
    place is sufficient (SD $< 0.3$ seats).

  \item Note that single-district states have point estimates without
    seed uncertainty (a single partition is defined by contiguity and
    population balance, not by METIS optimization).
\end{itemize}
